\chapter{Formal specification}

  \section{Type system}

 In this section, $\Rightarrow$ indicates the \texttt{synth}
 (inference) part of bidirectional typechecking while $\Leftarrow$
 indicates the \texttt{check} (verification) part. Since it trivially
 holds that $(x \underset{\textsc{tc}}{\Rightarrow} \tau) \Longrightarrow
   (x \underset{\textsc{tc}}{\Leftarrow} \tau)$, $\Leftrightarrow$ is simply written
   as $\Rightarrow$.
  
 Let $\iota \in \{ \texttt{Char}, \texttt{Bool}, \texttt{Int},
 \texttt{Float} \}$ and $c_\iota$ be a value of type $\iota$.
 \begin{equation*}
   \infer[\textsc{Lit}]{\Gamma \vdash c_\iota \Rightarrow \iota}{}\quad
   \infer[\textsc{Var}]
         {\Gamma \vdash x \Rightarrow \tau}
         {x: \tau \in \Gamma}
 \end{equation*}

 \begin{equation*}
   \infer[\textsc{Tup}]
         {\Gamma \vdash (e1, \ldots, e_n) \Rightarrow (\tau_1
           \times \ldots \times tau_n)}
         {\Gamma \vdash e_1 \Leftarrow \tau_1
           & \ldots
           & \Gamma \vdash e_n \Leftarrow \tau_n}
 \end{equation*}

 \begin{equation*}
   \infer[\textsc{Abs}]
         {\Gamma \vdash \texttt{fun } x: \sigma \texttt{ = } e
           \Rightarrow \sigma \rightarrow \tau}
         {\Gamma, x: \sigma \vdash e \Leftarrow \tau}
         \quad
         \infer[\textsc{App}]
         {\Gamma \vdash f\;a \Rightarrow \tau}
         {\Gamma \vdash a \Leftarrow \sigma
         & \Gamma \vdash f \Leftarrow \sigma \rightarrow \tau}
 \end{equation*}

 \begin{equation*}
   \infer[\textsc{PolyAbs}]
         {\Gamma \vdash \texttt{fun } x: \sigma \texttt{ = } e
         \Rightarrow \forall \alpha.\sigma \rightarrow \tau }
         {\begin{array}{c}
             \Gamma, \alpha \textrm{ type } \vdash \sigma \textrm{ type}\quad
             \Gamma, \alpha \textrm{ type } \vdash \tau \textrm{ type}
             \\
             \Gamma, \alpha \textrm{ type}, x : \sigma \vdash e
             \Leftarrow \tau
           \end{array}
         }
         \quad
         \infer[\textsc{PolyApp}]
         {\Gamma \vdash f\;a \Rightarrow \tau[\sigma/\alpha]}
         {\Gamma \vdash a \Leftarrow \phi(\sigma)
           &\Gamma \vdash f \Leftarrow \forall \alpha.\phi(\alpha) \rightarrow
           \tau}
 \end{equation*}

 \begin{equation*}
   \infer[\textsc{ADTConst}]
         {\Gamma \vdash k\;a_1 \ldots a_n \Rightarrow \tau\;\tau_1\;\ldots\;\tau_m}
         {\begin{array}{c}
             \Psi \vdash k \Leftarrow \forall \alpha_1, \ldots,
             \alpha_m.\sigma_1 \rightarrow \ldots \rightarrow
             \sigma_n \rightarrow \tau\;\alpha_1\;\ldots\;\alpha_m \in
             \mathcal{K}(\tau) \\
             \Gamma \vdash a_1 \Leftarrow
             \sigma_1[\tau_1/\alpha_1, \ldots, \tau_m/\alpha_m]\quad
             \ldots\quad
             \Gamma \vdash a_n \Leftarrow
             \sigma_n[\tau_1/\alpha_1, \ldots, \tau_m/\alpha_m] \\
           \end{array}
         }
 \end{equation*}

 \begin{equation*}
   \infer[\textsc{Let}]
         {\Gamma \vdash \texttt{let } x: \sigma \texttt{ = } y \texttt{ in } z
         \Rightarrow \tau}
         {\Gamma \vdash y \Leftarrow \sigma
         & \Gamma, x: \sigma \vdash z \Leftarrow \tau}
 \end{equation*}

 \begin{equation*}
   \infer[\textsc{If}]
         {\Gamma \vdash \texttt{if } e_i \texttt{ then } e_t \texttt{
             else } e_e \Rightarrow \tau}
         {\Gamma \vdash e_i \Leftarrow \texttt{Bool}
           & \Gamma \vdash e_t \Leftarrow \tau
           & \Gamma \vdash e_e \Leftarrow \tau}
 \end{equation*}

 \begin{equation*}
   \infer[\textsc{P\_}]
         {\Gamma \vdash \texttt{\_} \dashv \emptyset}
         {}
         \quad
         \infer[\textsc{Plit}]
         {\Gamma \vdash c_\iota \Rightarrow \iota \dashv \emptyset}
         {}
         \quad
         \infer[\textsc{Pbind}]
         {\Gamma \vdash x \Rightarrow \tau \dashv x \Leftarrow \tau}
         {}
 \end{equation*}

 \begin{equation*}
   \infer[\textsc{Pconst}]
            {\Gamma \vdash k\;p_1 \ldots p_n \Rightarrow
              \tau\;\tau_1\;\ldots\;\tau_m \dashv \Delta_1, \ldots,
              \Delta_n}
         {\begin{array}{c}
             \Psi \vdash k \Leftarrow \forall \alpha_1, \ldots,
             \alpha_m.\sigma_1 \rightarrow \ldots \rightarrow
             \sigma_n \rightarrow \tau\;\alpha_1\;\ldots\;\alpha_m \in
             \mathcal{K}(\tau) \\
             \Gamma \vdash p_1 \Leftarrow
             \sigma_1[\tau_1/\alpha_1, \ldots, \tau_m/\alpha_m] \dashv
             \Delta_1 \quad
             \ldots\quad
             \Gamma \vdash a_n \Leftarrow
             \sigma_n[\tau_1/\alpha_1, \ldots, \tau_m/\alpha_m] \dashv
             \Delta_n \\
           \end{array}
         }
 \end{equation*}

 \begin{equation*}
   \infer[\textsc{Ptup}]
         {\Gamma \vdash (p_1, \ldots, p_n) \Rightarrow (\sigma_1
           \times \ldots \times \sigma_n) \dashv \Delta_1, \ldots,
           \Delta_n}
         {\Gamma \vdash p_1 \Leftarrow \sigma_1 \dashv \Delta_1
           & \ldots
           & \Gamma \vdash p_n \Leftarrow \sigma_n \dashv \Delta_n}
 \end{equation*}

 \begin{equation*}
   \infer[\textsc{Por}]
         {\Gamma \vdash p \texttt{ | } q \Rightarrow \tau \dashv
           \Delta_p \cap \Delta_q}
         {\Gamma \vdash p \Leftarrow \tau \dashv \Delta_p
           & \Gamma \vdash q \Leftarrow \tau \dashv \Delta_q}
 \end{equation*}

 \begin{equation*}
   \infer[\textsc{Pat}]
         {\Gamma \vdash x\texttt{@}p \Rightarrow \tau \dashv \Delta, x
         \Leftarrow \tau}
         {\Gamma \vdash p \Leftarrow \tau \dashv \Delta}
 \end{equation*}

 \begin{equation*}
   \infer[\textsc{Match}]
         {\Gamma \vdash \texttt{match } e \texttt{ \{ } p_1 \texttt{
             if } c_1 \texttt{ -> } e_1, \ldots, p_n \texttt{ if } c_n
         \texttt{ -> } e_n \texttt{ \} } \Rightarrow \tau}
         {\begin{array}{c}
             \Gamma \vdash e \Leftarrow \sigma \\
             \Gamma \vdash p_1 \Leftarrow \sigma \dashv \Delta_1, e
             \Leftarrow \sigma\quad
             \ldots\quad
             \Gamma \vdash p_n \Leftarrow \sigma \dashv \Delta_n, e
             \Leftarrow \sigma \\
             \Gamma, \Delta_1 \vdash e_1 \Leftarrow \tau\quad
             \ldots\quad
             \Gamma, \Delta_n \vdash e_n \Leftarrow \tau\quad \\
             \Gamma, \Delta_1 \vdash c_1 \Leftarrow \texttt{Bool}\quad
             \ldots\quad
             \Gamma, \Delta_n \vdash c_n \Leftarrow \texttt{Bool} \\
           \end{array}
         }
 \end{equation*}

 \begin{multicols*}{2}
   In inference rules, $\Gamma$ denotes the general context whereas
   $\Psi$ specifically denotes the context defined by algebraic data
   types; $\Psi \subset \Gamma$.

   $\forall$- (polymorphic) types may only be constructed from arrow
   (function) types as in rule \textsc{PolyAbs}. The $\phi$ function
   in rule \textsc{PolyApp} can be understood as a purely syntactic
   rule for constructing types. It is used in this rule to allow for
   functions that do not directly take $\alpha$ as a parameter. For
   instance, the function \texttt{mapMaybe} is of type $\forall
   \alpha.\phi(\alpha) \rightarrow \forall \beta.(\alpha \rightarrow
   \beta) \rightarrow \phi(\beta)$, where $\phi(x) = \texttt{Maybe }
   x$. As can be seen in the previous example, whichever argument of the
   arrow type quantified using $\forall$ may itself contain other
   $\forall$ quantifiers.

   Data constructors are virtually functions, for the sake of
   consistency with the rules enunciated in the previous
   paragraph. The only difference is that data constructors are not
   curried. There are thus no nullary constructors, instead they are
   emulated as holding \texttt{()}. Thus, \texttt{Maybe} is defined as
   \texttt{data Maybe a = Just a | Nothing ()}.

   In pattern rules, a declaration such as $\Gamma \vdash x \Leftarrow \tau
   \dashv \Delta$ means that the pattern $x$ can be matched against
   type $\tau$ and that it provides context $\Delta$ (which is later
   used both by the guard $c_i$ and the expression $e_i$ in the
   \textsc{Match} rule).
 \end{multicols*}
